\subsection{Complete point lists in the rank-zero twist criteria}

\begin{statement}{A.33}{Elementary torsion bounds for the two CM families}{Proved}
For $b,c\in\Q^\times$, the rational torsion of
$E_b:y^2=x^3+b$ has order dividing six, and that of
$F_c:y^2=x^3+cx$ has order dividing four. The torsion points with both
coordinates nonzero on $E_b$ are precisely
\[
\begin{array}{c|c}
b & (x,y)\\\hline
-432r^6 &(12r^2,\ \pm36r^3)\\
r^6 &(2r^2,\ \pm3r^3)
\end{array}
\quad(r\in\Q^\times),
\tag{A33.1}
\]
and the torsion points on $F_c$ with $y\ne0$ are precisely
$c=4r^4$, $(x,y)=(2r^2,\pm4r^3)$.
The other possible finite torsion points of $E_b$ have $x=0$ when
$b$ is a square, or $y=0$ when $b$ is a cube.
\end{statement}
\begin{proof}
At a good prime $p\equiv2\pmod3$, the cube map on $\F_p$ is a
bijection, so summing the quadratic character of $x^3+b$ gives
$\#E_b(\F_p)=p+1$. At a good prime $p\equiv3\pmod4$, pairing $x$
and $-x$ gives $\#F_c(\F_p)=p+1$. Prime-to-$p$ rational torsion
injects under good reduction. Dirichlet's theorem, applied after
excluding the finitely many bad primes, gives the stated divisibility
bounds: for $E_b$ choose $p\equiv5\pmod{12}$ to bound the two-part,
$p\equiv2\pmod9$ to bound the three-part, and for any proposed
prime divisor $r>3$ choose $p\not\equiv-1\pmod r$ while
$p\equiv2\pmod3$. For $F_c$ use $p\equiv3\pmod8$ and exclude
each odd proposed divisor in the same way. All congruences used are
compatible reduced residue classes. This argument concerns arbitrary
rational coefficients, not only the later finite support census.

For $E_b$, two-torsion is $(-\sqrt[3]b,0)$ when rational. The
three-division equation is $3x(x^3+4b)=0$. If $x=0$ we get the
stated square case. Otherwise $x^3=-4b$ and $y^2=-3b$; putting
$r=y/(3x)$ gives $x=12r^2$, $y=36r^3$, and $b=-432r^6$.
This latter coefficient is not a rational cube, so this case has
no rational two-torsion. Order six therefore requires $b$ to be
both a square and a cube, hence $b=r^6$. Adding its two- and
three-torsion points gives $(2r^2,\pm3r^3)$. These lists exhaust
all possibilities because the torsion order divides six.

For $F_c$ the point $(0,0)$ has order two. A rational point of order
four must double to a rational point of order two. If all three
nonzero two-torsion points are rational, the order bound four leaves
no point of order four. Otherwise its double is $(0,0)$, and the
duplication formula gives $x^2=c$, $y^2=2cx$. Hence $c=4r^4$,
$x=2r^2$, $y=\pm4r^3$. Direct substitution and duplication verify
these points. In the finite support census, the two bounds also have
the particularly short certificates
$\gcd(\#E_b(\F_{11}),\#E_b(\F_{17}))=6$ and
$\gcd(\#F_c(\F_{11}),\#F_c(\F_{19}))=4$.
\end{proof}

\begin{statement}{A.34}{All ten elliptic factors of a diagonal sextic}{Proved}
Let $X_{\alpha,\beta,6}$ be the smooth projective curve
$\alpha X^6+\beta Y^6=Z^6$. Its Jacobian is isogenous over $\Q$
to the product of ten elliptic curves $y^2=x^3+B$ with coefficients
\[
\begin{array}{c|l}
\text{type}&B\text{ modulo rational sixth powers}\\\hline
\text{six triangle maps}&
 \alpha^2\beta^3,\ \alpha^3\beta^2,\ \alpha^2\beta,\
 \alpha\beta^2,\ -\alpha^3\beta,\ -\alpha\beta^3\\
\text{three cyclic maps}&
 -4\alpha\beta,\ 4\alpha/\beta^2,\ 4\beta/\alpha^2\\
\text{cubic map}&-432(\alpha\beta)^2.
\end{array}
\tag{A34.1}
\]
Repeated isomorphism classes in this list retain their multiplicities.
Each factor has geometric endomorphism ring $\Z[\zeta_3]$.
Nonvanishing of $L(E,1)$ for any one factor suffices for A.24's
point-or-cusp conclusion on the sextic open curve.
\end{statement}
\begin{proof}
Put $(a_0,a_1,a_2)=(\alpha,\beta,-1)$ and
$(X_0,X_1,X_2)=(X,Y,Z)$. For every permutation $(i,j,k)$, the functions
\[
x=-a_i a_j(X_i/X_k)^2,\qquad
y=a_i a_j^2(X_j/X_k)^3
\]
satisfy $y^2=x^3-a_i^2a_j^3a_k$. This gives the six triangle maps
over $\Q$. The three remaining cyclic maps arise as follows. In the
chart $X_k=1$ put
\[
t=-a_iX_i^6/a_k,\qquad w=X_iX_j,\qquad v=2t-1.
\]
Then $v^2=1-4a_i a_jw^6/a_k^2$. Consequently
$x=w^{-2}$, $y=vw^{-3}$ give an elliptic target with
$B=-4a_i a_j/a_k^2$, one for each unordered pair $\{i,j\}$.
The cubic map and its Jacobian are A.26 and A.27. Every displayed
nonconstant map extends to the smooth projective models. For the
cubic torsor target, compose with its degree-three divisor map to its
Jacobian, as in A.24; this remains nonconstant and defined over $\Q$.

We check that these ten maps account for the whole Jacobian, rather
than merely ten possible targets. Over $\overline\Q$, in the chart
$Z=1$, a basis of regular differentials is indexed by the ten pairs
$(a,b)$ with $a,b\ge1$ and $a+b<6$, represented up to nonzero scalar by
$X^{a-1}Y^{b-6}\,dX$. The six triangle maps pull back to the characters
given by the six permutations of $(1,2,3)$, taking the first two
entries. The cyclic maps pull back to the three permutations of
$(1,1,4)$. For example the map attached to $\{X,Y\}$ gives
\[
\frac{dx}{2y}
=-\frac{d(XY)}{2\alpha X^6-1}
=\frac{dX}{\beta Y^5},
\]
using $\alpha X^6+\beta Y^6=1$ and its differential.
The cubic map gives character $(2,2)$, from the triple $(2,2,2)$.
These ten distinct lines are exactly the displayed differential basis.
The induced homomorphism from the Jacobian to the product of the ten
targets is therefore an isomorphism on tangent spaces. In characteristic
zero a homomorphism of abelian varieties of equal dimension with this
property is an isogeny. Finally each elliptic equation has $j=0$ and the
automorphism $(x,y)\mapsto(\zeta_3x,y)$; maximality of
$\Z[\zeta_3]$ proves the endomorphism assertion. Apply A.24 and Rubin's
rank-zero theorem to any one factor. The code \texttt{root\_exact.py}
checks all nine rational-map identities and the differential indexing.
\end{proof}

\begin{statement}{A.35}{Universal finite lists of ordinary point sections}{Proved}
Suppose a diagonal twist satisfies one of the rank-zero criteria of
A.24--A.27 or A.34. Its ordinary rational points above the tripod are
obtained from the following complete lists of possible tripod values:
\[
\begin{aligned}
\mathcal T_3=\mathcal T_4&=\{-1,\ 2,\ 1/2\},\\
\mathcal T_6&=\mathcal T_3\ \cup\ \mathfrak S_3\cdot(-8)
                         \ \cup\ \mathfrak S_3\cdot4.
\end{aligned}
\tag{A35.1}
\]
Here the action is generated by $t\mapsto1-t$ and $t\mapsto1/t$;
$\mathcal T_6$ has fifteen distinct elements. A value $t$ in the
relevant list occurs precisely when
\[
[t]=[\alpha],\qquad[1-t]=[\beta]
          \quad\text{in }\Q^\times/\Q^{\times n}.
\tag{A35.2}
\]
Every point upstairs is then
\[
[X:Y:Z]=\left[\sqrt[n]{t/\alpha}:
                       \sqrt[n]{(1-t)/\beta}:1\right]
\tag{A35.3}
\]
with all rational choices of roots. The quartic classes with an ordinary
point are $(-1,2),(2,-1),(8,8)$, corresponding to $t=-1,2,1/2$.
For cubic classes the analogous representatives are $(1,2),(2,1),(4,4)$.
This is a complete list of point sections under the criteria, and is
independent of any height-bounded search.
\end{statement}
\begin{proof}
The cubic case is also proved in the census below. Here is a geometric
derivation. If the compact diagonal cubic has a rational point $P$, it
is a torsor trivialized by $P$. The rational hyperplane divisor $H$ defines
the rational point $3[P]-H$ on its Jacobian $E_d$, $d=[\alpha\beta]$.
Rank zero and A.33 imply that $E_d(\Q)$ is trivial unless $d=1$ or $2$,
with torsion $C_3$ or $C_2$ in those two cases. In the trivial case
$3[P]\sim H$, so $P$ is a flex; the Hessian of a diagonal cubic is a
nonzero constant times $XYZ$, making every such point a coordinate cusp.

If $d=1$, a nontrivial coefficient class has at some prime the three
coefficient valuations congruent to $0,1,2$ modulo three. In any
putative local point the three nonzero terms would then have distinct
valuation classes, and the smallest valuation could not cancel.
The only soluble class is $(1,1)$, whose rational points are its three
cusps by the classical cubic calculation. If $d=2$, the same local
argument away from two forces both coefficients to be supported only
at two, and gives the three classes $(1,2),(2,1),(4,4)$.
Each has an explicit cusp and an explicit ordinary point, respectively
$[-1:1:1]$, $[1:-1:1]$, and $[1:1:2]$. Its Jacobian has exactly two
rational points, so these exhaust the rational points of the torsor.
Their tripod values are $-1,2,1/2$.

For the quartic, take whichever genus-one quotient in A.25 has rank zero,
and write it as $v^2=A u^4+B$. It maps over $\Q$ to
\[
F_{AB}:y^2=x^3+ABx,\qquad x=A u^2,\quad y=Auv.
\]
This target is two-isogenous to its Jacobian, hence also has rank zero.
An ordinary original point has $u\ne0$ and $v$ a nonzero rational
square. Thus $y\ne0$, and A.33 forces
$AB=4r^4$, $Au^2=2r^2$, $Auv=\pm4r^3$.
Eliminating $r$ gives
\[
A u^4=B,\qquad v^2=2B.
\]
Since $v$ is a rational square, $B$ is eight times a rational fourth
power, and the same holds for $A$. In the first quotient
$(A,B)=(\alpha,\beta)$ this gives class $(8,8)$ and $t=1/2$.
In the other two quotients it gives $(-1,2)$ and $(2,-1)$, with
$t=-1$ and $2$. Direct substitution supplies all the roots in (A35.3).

For a sextic triangle quotient, first use the ordering
$\alpha X^6+\beta Y^6=1$ and
$x=-\alpha\beta X^2$, $y=\alpha\beta^2Y^3$ on
$E_{\alpha^2\beta^3}$. At an ordinary point both coordinates are
nonzero, so A.33 gives two cases. If $\alpha^2\beta^3=r^6$,
the torsion point has $(x,y)=(2r^2,\pm3r^3)$. Cubing the $x$ equation
and dividing by $r^6$ yields $-\alpha X^6=8$, hence $t=-8$.
If $\alpha^2\beta^3=-432r^6$, the torsion point has
$(x,y)=(12r^2,\pm36r^3)$, giving $t=4$ in the same way.
Permuting the three coordinates yields exactly the two six-element
orbits in (A35.1).

For an extra cyclic target, $x=w^{-2}$ is a nonzero rational square.
The nonzero $x$-coordinates in (A33.1) are respectively three and two
times rational squares, so neither can be a square. Thus the elliptic
image must be two-torsion, with $y=0$. The formula $y=(2t-1)w^{-3}$
forces $t=1/2$. The other coordinate orderings give its three-element
orbit. Finally, if it is the cubic target that has rank zero, apply
the already proved cubic case to the squaring map. This proves necessity
of the finite list in every criterion. Conversely (A35.2) gives actual
rational roots, so (A35.3) proves sufficiency and completeness.
\end{proof}

\begin{statement}{A.36}{Complete Selmer section sets under the CM criteria}{Proved}
For every diagonal cubic, quartic or sextic twist satisfying at least
one of the stated analytic rank-zero criteria, its Selmer section set
is exactly the disjoint union of the ordinary point sections (A35.3)
and the full packets at its rational coordinate cusps. The latter are
obtained by solving the corresponding two-term equation after setting
one coordinate to zero. On the tripod, its Kummer fibre is the union
of the point classes (A35.2) and the three packet cosets in A.23.
Under A.20 the same descriptions hold for all full sections.
In particular Lead 2's geometric conclusion is true, but the additional
finiteness assertion in tasks A.5--A.6 is false whenever one of those
packet cosets is present.
\end{statement}
\begin{proof}
Rubin's theorem and A.24 prove that no other Selmer sections occur.
A.35 enumerates every ordinary rational point; A.23 enumerates every
downstairs cusp fibre. Rational tangents at each rational cusp identify
the upstairs packet with $H^1(\Q,\wh\Z(1))$. The full-section statement
uses precisely A.20 and A.21. The explicit infinitely many parameter
values in A.23 prove failure of finiteness. There is no quotient by
these parameters implicit in the word ``packet''.
\end{proof}
